\begin{figure}[htbp]
\centering
\begin{tikzpicture}[
  >=Latex, font=\footnotesize,
  bx/.style={draw, rounded corners=1.5pt, minimum height=5mm,
             inner xsep=3.5pt, inner ysep=1.5pt},
  lb/.style={inner sep=1.5pt},
  tag/.style={fill=white, inner sep=1.5pt, font=\scriptsize}]
\node[lb] (st) at (0,0) {$s_t$};
\node[bx] (mt) at (1.15,0) {$\mu_{\theta,t}$};
\node[bx] (pt) at (2.80,0) {$\bar P_t$};
\node[bx] (ft) at (4.35,0) {$f_t$};
\node[lb] (sn) at (5.70,0) {$s_{t+1}$};
\node[lb] (qt) at (2.80,-1.20) {$b_t(z_t)+B_ts_t$};
\node[lb] (xt) at (4.35,0.80) {$\xi_{t+1}$};

\draw[->] (st) -- (mt);
\draw[->] (mt) -- node[above,lb] {$u_t$} (pt);
\draw[->] (pt) -- node[above,lb] {$x_t$} (ft);
\draw[->] (ft) -- (sn);
\draw[->] (qt) -- (pt);
\draw[->] (xt) -- (ft);

\node[bx] (mn) at (8.65,1.05) {$\mu_{\theta,t+1}$};
\node[lb] (qn) at (8.65,-1.20) {$b_{t+1}(z_{t+1})+B_{t+1}s_{t+1}$};
\node[bx] (pn) at (12.05,0) {$\bar P_{t+1}$};
\draw[->,draw=blue!60!black,thick] (sn) to[out=55,in=180]
  node[tag,text=blue!60!black,pos=0.55] {proposal--state} (mn.west);
\draw[->,draw=orange!85!black,thick] (sn) to[out=-55,in=155]
  node[tag,text=orange!85!black,pos=0.55] {moving recourse} (qn.north west);
\draw[->] (mn.east) to[out=0,in=115]
  node[lb,above right,pos=0.15] {$u_{t+1}$} (pn.north);
\draw[->] (qn.east) to[out=0,in=-115] (pn.south);
\draw[->,dashed] (pn) -- ++(0.8,0) node[right,lb] {$\cdots$};
\end{tikzpicture}
\ReviewBegin%
\caption{Two consecutive stages of the projected policy. The proposal map
$\mu_{\theta,t}$ produces $u_t$, the projection $\bar P_t$ produces a feasible
decision $x_t$, and the dynamics $f_t$ produce the next state $s_{t+1}$.
This state affects the next projection through both the proposal (blue)
and the constraint bound $b_{t+1}(z_{t+1})+B_{t+1}s_{t+1}$ (orange).
Holding the bound fixed during differentiation removes the orange path
while leaving the forward simulation unchanged. Here $z_t$ summarizes the
observed exogenous history; other inputs to the proposal and dynamics are
omitted for clarity.}
\label{fig:reservoir-intuition}
\ReviewEnd%
\end{figure}
